\begin{saresult}{A.131}{Proved under stated hypotheses}{Effective countable models and complete recursive presentations for PL and semialgebraic concordances}
Let $K$ be a nonempty finite abstract simplicial complex whose realization
is a closed PL manifold. The ordered finite list of simplices of $K$ is
all the input in the following construction. There are explicit
countable simplicial groups
\[
 G^{\mathrm{PL}}(K),\qquad G^{\mathrm{SA}}(K)
                                                        \tag{A.131.1}
\]
with the following effective properties. Their simplices are finite
integer strings; membership, equality, and degeneracy are decidable.
Multiplication, inversion, faces and degeneracies are given by total
algorithms. The construction also gives a computable simplicial
homomorphism $G^{\mathrm{PL}}(K)\to G^{\mathrm{SA}}(K)$.
For every real closed field $R$, coefficient extension gives weak
equivalences
\[
 G^{\mathrm{PL}}(K)\longrightarrow P_{\mathrm{PL},R}(|K|_R),
 \qquad
 G^{\mathrm{SA}}(K)\longrightarrow P_{\mathrm{SA},R}(|K|_R),
                                                        \tag{A.131.2}
\]
where the source is realized over the real algebraic numbers.
All families are homeomorphisms of the total product with the
parameter simplex, preserve that entire parameter, fix the bottom
pointwise, and preserve the free top setwise.

For each $r\ge0$, the algorithm in the proof outputs a complete
countable presentation $\Pi_r^{\mathrm{PL}}(K)$, respectively
$\Pi_r^{\mathrm{SA}}(K)$, for the entire homotopy group
\[
 \pi_rP_{\mathrm{PL},R}(|K|_R),\qquad
 \pi_rP_{\mathrm{SA},R}(|K|_R).
                                                        \tag{A.131.3}
\]
For $r\ge1$ these groups are based at the identity; translation
identifies the corresponding groups on any other component.
Here ``outputs'' means that one uniform algorithm enumerates all
specified generators and relators. For $r\ge1$ it may instead
enumerate an integer relation matrix for the abelian group.
For an arbitrary fixed o-minimal expansion, $\Pi_0^{\mathrm{PL}}(K)$
also presents the stipulated definable component group. If every
component of $|K|$ has dimension at least five, (A.131.1) gives a
weak model for the real TOP concordance space as well, and
$\Pi_r^{\mathrm{PL}}(K)$ presents its homotopy groups for all $r$.

For the empty complex and empty base, take the constant one-element
simplicial group; every component and homotopy group is trivial.

These are complete effective countable presentations and whole
simplicial models. No terminating word problem, finite presentation,
finite invariant-factor description, or efficient complexity bound
is asserted. No higher weak equivalence from PL to SA or to an
arbitrary o-minimal expansion is inferred.
\end{saresult}
\begin{proof}
\emph{The finite coordinate model and its arithmetic.}
Order the $v$ vertices of $K$ and realize them as the standard basis
vectors of $\mathbb Q^v$. Put
\[
 B=|K|,\quad Z=B\times[0,1],\quad A=B\times\{0\},\quad
 T=B\times\{1\},\quad D_q=Z\times\Delta^q.
\]
Membership in $B$ is the finite disjunction, over the maximal
simplices $\sigma$ of $K$, of
\[
 x_i\ge0,\qquad \sum_i x_i=1,\qquad x_i=0\ (i\notin\sigma).
\]
This is a fixed rational linear formula. The other displayed sets have
rational linear formulas too. None of the algorithms needs to
recognize whether the promised finite complex is a manifold.

Write $F=\mathbb Q^{\mathrm{rc}}$ for the ordered real algebraic
numbers. One may encode an element by a nonzero integer polynomial,
a rational open interval with exactly one root, and that root.
The final phrase refers to its unique value, not additional input:
the finite input is the polynomial and the two rational endpoints.
Sturm root counting decides whether this code is valid. Every real
algebraic number has such a code, and the set of valid codes is
effectively enumerable. Arithmetic, sign and equality can be decided
by root isolation, or by the following uniform reduction to rational
quantifier elimination. Replace each coded number by a fresh variable,
constrain it by its polynomial and isolating interval, and ask the
corresponding polynomial sentence. Each such variable has exactly one
possible value, so this replacement preserves truth. An output code
for an arithmetic operation is found by enumerating root codes and
deciding equality to the desired expression; a code exists, so the
search terminates. Duplicate codes for a number cause no difficulty.

The exact quantifier-elimination input is
\cite[online edition dated 26 June 2016, Theorem~2.94, p.~80;
Algorithm~14.16, pp.~586--587; Algorithm~14.21, pp.~591--592,
and Theorem~14.22, pp.~592--593]{BasuPollackRoy2016Online}.
The same edition states Sturm's theorem as Theorem~2.62, p.~62.
Thus the arithmetic used below is an algorithm on finite integers,
not arithmetic with an oracle for arbitrary real constants.
For linear formulas over this effective ordered field there is also
an elementary elimination algorithm: put a Boolean formula into
disjunctive normal form; substitute an equality with nonzero
coefficient when present; otherwise compare each lower bound with
each upper bound for the variable being eliminated. Retain a strict
comparison if either original bound was strict. A finite family of
one-sided bounds is satisfiable when there is no opposite bound,
and constant inequalities are decided by their signs. Negating
converts universal to existential quantifiers. Iterating eliminates
all variables. This proves the linear decision claim used below.

\emph{PL graphs and their finite validity test.}
In degree $q$ use ambient coordinates $(z,z',t)$, where
$z,z'\in F^{v+1}$ and $t\in F^{q+1}$. In particular the source
and target share the \emph{same} parameter variable $t$.
A raw code is a finite list of closed simplices with vertices in $F$
in this ambient space. Affine independence is decided by Gaussian
elimination over $F$; reject a purported simplex that is not
independent, and allow zero-dimensional simplices. The list need not
already form a simplicial complex. Its union $\Gamma$ has the formula
\[
 \bigvee_{\sigma}
 \left\{w=\sum_{j=0}^{a_\sigma}\lambda_jv_{\sigma j},\quad
            \lambda_j\ge0,\quad\sum_j\lambda_j=1\right\}.
                                                        \tag{A.131.4}
\]
The braces abbreviate existential quantification of the finitely
many $\lambda_j$. All coefficients are constants in $F$, so
(A.131.4) is a linear formula.

Call the raw code valid precisely when these five finite sentences
hold:
\begin{enumerate}
\item $\Gamma\subset Z\times Z\times\Delta^q$;
\item for every $(z,t)\in D_q$ there exists exactly one $z'$ with
$(z,z',t)\in\Gamma$;
\item for every $(z',t)\in D_q$ there exists exactly one $z$ with
$(z,z',t)\in\Gamma$;
\item whenever $(z,z',t)\in\Gamma$ and $z\in A$, one has $z'=z$;
\item whenever $(z,z',t)\in\Gamma$, one has $z\in T$ if and only if
$z'\in T$.
\end{enumerate}
Uniqueness is the formula with two candidate outputs implying their
equality. Every sentence is linear and therefore decidable by the
elimination just given. Equality of two graph codes is the linear
sentence asserting their membership predicates equivalent at every
point, so it is decidable too.

We prove that this test is sufficient, including continuity. The
projection of $\Gamma$ to $(z,t)$ is bijective by the first three
conditions. On each simplex $\sigma$ it is affine and injective.
An affine map injective on a simplex is injective on its affine hull:
a kernel direction would give two distinct nearby points in the
relative interior with the same image. Its inverse on its image is
therefore affine. It follows that the function determined by
$\Gamma$ is affine on each of the finitely many projected simplices.
Those simplices are closed in $D_q$, and their affine formulas agree
on their overlaps because the graph is single valued. The finite
closed-cover gluing lemma proves continuity over $F$. Applying the
same argument to the other projection proves inverse continuity.

For completeness, a finite polyhedral subdivision compatible with
all the projected simplices is obtained by cutting with their affine-hull equations and ambient
extensions of all their face supporting hyperplanes, taking every
resulting face in every relevant affine span, and
triangulating those cells inductively by coning from an interior
point. A bounded nonempty polyhedron over $F$ has vertices over $F$,
by Gaussian elimination; the average of the vertices lies in its
relative interior. Thus all these subdivisions have vertices in $F$.
The affine formulas are compatible on this subdivision. The function
and its inverse are jointly PL, and the final two sentences impose
exactly the stipulated boundary conditions. Here PL means finite piecewise linear, as in A.115. Conversely, a jointly PL
homeomorphism on the finite polyhedron $D_q$ is affine on a finite
subdivision; its graph on each simplex is a simplex of the prescribed
kind. Hence every stipulated PL family has a valid graph code.
This equivalence used neither a homeomorphism-recognition oracle nor
ordinary compactness of a non-Archimedean field.

\emph{Canonical representatives and total operations.}
Encode all finite syntax above by binary strings using fixed
self-delimiting encodings of signed integers, lists and root codes.
Order the strings first by length, then lexicographically. There
are only finitely many predecessors of a given string. The canonical
representative of a valid graph is its least valid raw code in that
order. Given a valid code $c$, check the finitely many strings at
most $c$ for validity and graph equality, and take the first match.
Thus canonicalization terminates, and the predicate saying that a
string is canonical is decidable. Define $G_q^{\mathrm{PL}}(K)$ to
be the set of canonical valid degree-$q$ codes.

Here are finite formulas for all operations. If $g$ and $h$ are
represented by graph predicates $\Gamma_g$ and $\Gamma_h$, put
\[
 \begin{split}
 \Gamma_{g^{-1}}(z,z',t)&=\Gamma_g(z',z,t),\\
 \Gamma_{gh}(z,z',t)&=\exists y\,
       (\Gamma_h(z,y,t)\wedge\Gamma_g(y,z',t)).
 \end{split}
                                                        \tag{A.131.5}
\]
The product convention is $g$ after $h$.
For an order-preserving map $\alpha:[m]\to[q]$, let
\[
 (\alpha_*s)_j=\sum_{\alpha(i)=j}s_i.
\]
The simplicial pullback has graph
\[
 \Gamma_{\alpha^*g}(z,z',s)=
     [\,(z,s)\in D_m,\ (z',s)\in D_m\,]
       \wedge\Gamma_g(z,z',\alpha_*s).
                                                        \tag{A.131.6}
\]
This covers degeneracies as well as faces. The identity graph is
$z'=z$ on $D_q$.

All these are linear formulas. They are graphs of jointly PL
homeomorphisms: inversion exchanges the affine subdivisions;
composition uses intersections of the image cells of $h$ with the
domain cells of $g$; pullback uses the inverse images of the cells
under the affine map $(z,s)\mapsto(z,\alpha_*s)$ and the formula
for the first coordinate of $g$. Their inverse maps are obtained by
the same construction. These are finite bounded polyhedral covers,
triangulable over $F$ by the construction above. Consequently each
operation formula has at least one valid raw graph code.

The following uniform search outputs its canonical code: enumerate
raw codes in the prescribed order, reject invalid codes, and test
whether their graph predicate equals the operation formula. Equality
is decidable and a successful code exists, so the search terminates.
This proves total computability of every operation without leaving
a triangulation subroutine unspecified. Canonical equality is actual
equality of graphs, so associativity, inverse laws, and all
simplicial identities hold strictly. We have constructed a simplicial
group literally isomorphic to the jointly PL concordance group over
$F$.

\emph{The semialgebraic construction.}
For $G^{\mathrm{SA}}$ replace a list of graph simplices by a finite
Boolean formula in polynomial equations and inequalities with coded
coefficients in $F$, in the same variables $(z,z',t)$. Use the same
five validity conditions. Add joint continuity of the graph function
and of its inverse. These are first-order sentences, not assertions
that each individual fibre map is continuous. For example continuity
of the function is the sentence
\[
 \begin{split}
 \forall (z,z',t)\in\Gamma\ \forall\epsilon>0\ \exists\delta>0\
 \forall (w,w',s)\in\Gamma:
 \quad &\|w-z\|^2+\|s-t\|^2<\delta^2\\
 &\Longrightarrow
 \|w'-z'\|^2+\|s-t\|^2<\epsilon^2.
 \end{split}
                                                        \tag{A.131.7}
\]
Here squared norms are finite sums of coordinate squares, so the
formula is polynomial. Exchange the two graph projections to impose
inverse continuity. These conditions express precisely that the
total-space bijection is a homeomorphism. Quantifier elimination
decides their truth and equality of graphs. It also eliminates the
quantifiers in (A.131.5) when composition is formed. The same
least-string procedure produces canonical codes and total
algorithms. It follows exactly as before that
$G^{\mathrm{SA}}(K)$ is a simplicial group, literally the stipulated
semialgebraic concordance group over $F$.

Formula (A.131.4) is already a semialgebraic graph formula. Canonicalize
it in the SA coding to obtain the indicated computable map
$G^{\mathrm{PL}}\to G^{\mathrm{SA}}$. The map respects every
operation because each operation was defined by its actual graph.
A.115 and A.116, applied to $F\subset R$, prove (A.131.2).
The unique ordered embedding of the real algebraic field into any
real closed field specifies those coefficient extensions. There is
no map from non-Archimedean points to ordinary real points in this
argument.

These are whole effective simplicial models, not merely lists of
candidate components. For example, a degree-$q$ canonical simplex
$a$ is degenerate exactly when
\[
 a=s_i d_i a\quad\hbox{for some }0\le i<q.
\]
The finite test is decidable. Thus one can enumerate the
nondegenerate cells of the geometric realization and compute each
of their face attachments. This specifies its entire countable CW
model, with all incidence maps, in addition to its simplicial group
structure.

\emph{All homotopy groups with every relation.}
Let $G$ denote either effective simplicial group. Put
\[
 Z_0=G_0,\qquad
 Z_r=\bigcap_{i=0}^{r}\ker(d_i:G_r\to G_{r-1})\quad(r\ge1),
 \qquad
 N_{r+1}=\bigcap_{i=1}^{r+1}\ker(d_i:G_{r+1}\to G_r).
                                                        \tag{A.131.8}
\]
Membership in these sets is decidable, since there are finitely many
computable face tests. The simplicial identities show
$d_0N_{r+1}\subset Z_r$. This is a subgroup, and it is normal in
$Z_r$: for $a\in Z_r$ and $u\in N_{r+1}$,
\[
 (s_0a)u(s_0a)^{-1}\in N_{r+1},\qquad
 d_0\big((s_0a)u(s_0a)^{-1}\big)=a(d_0u)a^{-1}.
                                                        \tag{A.131.9}
\]
The same argument works in degree zero with $Z_0=G_0$.

Every simplicial group is Kan, and its Moore quotient computes its
homotopy groups. The precise opened source is
\cite[Theorem~17.1, pp.~67--68; Proposition~17.2, p.~68;
the second Proposition~17.3, pp.~68--69; Proposition~17.4, p.~69;
Theorem~16.6, p.~65]{May1967Simplicial}.
The original prints two propositions numbered 17.3; the second is
the normal-subgroup statement. It uses the last face convention,
whereas (A.131.8) uses the first. Reversing the order on each
standard simplex exchanges these conventions. More directly, if
$z$ represents an elementary simplicial homotopy with
\[
 d_0z=x,\quad d_1z=y,\quad d_i z=1\quad(2\le i\le r+1),
 \qquad x,y\in Z_r,
\]
then $u=z(s_0y)^{-1}$ belongs to $N_{r+1}$ and
$d_0u=xy^{-1}$. Conversely $u\in N_{r+1}$ produces the same
homotopy from $u s_0y$ whenever $d_0u=xy^{-1}$.
For $r=0$ this is the corresponding edge calculation; the Kan
condition composes edge paths. The multiplication on these
homotopy classes is induced by the degreewise group product,
by Proposition~17.2. Hence there is an isomorphism of groups
\[
 \pi_r|G|\cong Z_r/d_0N_{r+1}\quad(r\ge0).
                                                        \tag{A.131.10}
\]
The realization comparison in Theorem~16.6 identifies the simplicial
and ordinary homotopy groups used here.

The promised presentation is completely specified as follows:
\[
 \Pi_r(G)=
 \left\langle x_a\ (a\in Z_r)\ \left|\
 x_1=1,\quad x_ax_b=x_{ab}\ (a,b\in Z_r),\quad
 x_{d_0u}=1\ (u\in N_{r+1})\right.\right\rangle.
                                                        \tag{A.131.11}
\]
Enumerate the canonical binary codes, use the finite tests in
(A.131.8), and dovetail their pairs to enumerate the displayed
relators. The product and face indices are computed by the already
specified total algorithms. Thus no word in (A.131.11) is indexed
by an unspecified isotopy or an undecidable predicate.
The first two families of relators present $Z_r$ via its full
multiplication table: mapping $x_a$ to $a$ and $a$ to $x_a$ gives
inverse homomorphisms. The last family takes the quotient by the
normal subgroup $d_0N_{r+1}$, which is precisely the subgroup in
(A.131.10). This proves that every relation, and only the required
relations, is included.

For $r\ge1$ the quotient is abelian, as also proved in the first
Proposition~17.3 on p.~68 of the same source. Therefore one may
instead enumerate the free abelian generators $e_a$ and rows
\[
 e_1,\qquad e_a+e_b-e_{ab},\qquad e_{d_0u}
                                                        \tag{A.131.12}
\]
of a countable integer relation matrix. Abelianization adds no
further relation to the already abelian group (A.131.11).
The subgroup $d_0N_{r+1}$ is recursively enumerable, since its
preimages have a decidable membership test, but its membership is
not claimed decidable. In particular none of this asserts an
algorithm that answers every isotopy question in finite time.

Finally A.113 identifies $\Pi_0^{\mathrm{PL}}(K)$ with the
definable component group in every fixed o-minimal expansion.
In the stated dimension-at-least-five range the same result gives
a PL--TOP weak equivalence over the real field, and hence the full
TOP model and all homotopy presentations asserted above.
This uses the comparison with the singular complex of the TOP
group, not a literal real-points map from an arbitrary field.

\emph{Finite verification and the exact meaning of the output.}
The attached independent program checks a rational triangulated
subclass of (A.131.4) for $B$ a point and $q=0,1,2$ using exact
rational arithmetic. It constructs the common subdivisions in
(A.131.5) and affine preimages in (A.131.6), tests both projection
triangulations, and checks nonconstant jointly PL families
\[
 h_t(x)=x+\tfrac12\min(x,1-x,t_0,\ldots,t_q)
 \quad(q=1,2).
\]
It checks inverse identities, subdivision equality, all faces of the
two-simplex example, and degeneracies of a nonidentity zero-simplex;
it rejects five defective graph codes. All thirty checks passed,
and an independent author rerun passed in $0.842377$ seconds.
A second author run, saved separately, passed in $0.761886$ seconds,
and the root reviewer independently passed all thirty in $0.778375$
seconds.
The program does not claim to implement general polynomial
quantifier elimination or certify the infinite Moore theorem.

The result is an algorithm from a finite complex to a countable
simplicial group with decidable simplices and computable structure,
and thence to complete recursively enumerable group presentations.
The real and arbitrary-field concordance spaces are identified with
that model by the proved coefficient comparisons. This gives more
than a name for a concordance group: every finite candidate and
every displayed relation is checked by terminating algebraic
procedures. It remains an exhaustive presentation, not an evaluation
into finite classical invariants. The separate question whether the
computable map $G^{\mathrm{PL}}\to G^{\mathrm{SA}}$ is a weak
equivalence is unchanged. For higher families in an arbitrary
fixed o-minimal expansion, no effective reduction to these algebraic
codes is asserted; that comparison remains a distinct gate.
\end{proof}
